\chapter{Conclusion}\label{ch:conclusion}

\section{Summary}\label{sec:conc:summary}

This thesis replaced the diffusion model of \ac{DPCC} with three flow-based objectives,
instantaneous-velocity matching, analytic average-velocity matching and consistency-interpolated
average-velocity matching, and set endpoint projection beside the per-step projection of \ac{DPCC},
keeping the temporal U-Net, the projection program, the data and the evaluation of the baseline. The
step budget thereby moved from training to inference. The comparison was made in three steps: on the
obstacle-avoidance benchmark of \ac{DPCC} with its setup unchanged, on a vision-conditioned alignment
task built for this thesis, and on a quadrotor benchmark of three scenes built for it as well.

On the benchmark of \ac{DPCC} the average-velocity models Pareto-dominate its diffusion model with its
own U-Net: the same success with constraint satisfaction, fewer control steps, and about a thirtieth
of the time per action, twelve times less than a baseline retrained at two denoising steps.
Instantaneous-velocity matching dominates the baseline as well, so the saving belongs first to a
flow-based sampler at a small budget and then to the average-velocity objective, which adds the fewest
control steps. On D3IL-aligning the dominance repeats where the task separates the models: analytic
average-velocity matching is the only model that moves the box in every context, it closes most of
the distance to the target, and it is the only objective that improves with the budget, while the
consistency-interpolated model, instantaneous-velocity matching and the baseline barely move the box.
Endpoint projection has no step to guide at the one-evaluation operating point of D3IL-avoiding and
nothing to win there; on D3IL-aligning, run at twenty evaluations, it is the better projector on the
constraint, keeping every context violation-free at the price of per-step projection.

The quadrotor scenes carry the framework to a plant that holds nothing by itself. UAV-pillars flies
the operating points of D3IL-avoiding and keeps their result on the declared constraints, and shows
what the plant changes: the vehicle tracks more closely than the arm, and a collision with a pillar the
constraints leave out ends the flight. UAV-corridor shows that the projection scheme of \ac{DPCC}
carries over to three dimensions, against constraints that leave the horizontal plane and that no
demonstration satisfies, and that on straight-line demonstrations the generative model does not
matter, the budget does; there the two projectors trade against each other without one dominating.
UAV-s-curve shows that the controller matters: the same plans succeed or fail by the controller that
flies them, and the cascaded geometric controller used throughout is the better of the two, at a
twenty-fourth of the cost.

With human demonstrations the average-velocity models lead, on the harder tasks the
model and the projection matter together and endpoint projection helps where the task needs a large
budget; with simple generated demonstrations instantaneous-velocity matching is enough and the
objectives are equivalent; and on the quadrotor the plant and its controller matter as much as the
planner.

\section{Answers to the Research Questions}\label{sec:conc:rqs}

\begin{description}
  \item[RQ1 --- Generative model.] Yes, with a margin. With the backbone, the data and the evaluation
    held fixed, the flow-based models reach the diffusion baseline's success with constraint
    satisfaction at one network evaluation, where the diffusion model, run away from its trained
    budget, loses it: on D3IL-avoiding consistency-interpolated average-velocity matching holds
    $1.000$ in $59.2$ control steps at $18.1$\,ms per action against the baseline's $1.000$ in $70.1$
    at $553.4$, and instantaneous-velocity matching $1.000$ in $67.0$ at $17.3$. Among the objectives
    the order, wherever the task separates them, is the average-velocity models, then
    instantaneous-velocity matching, then the diffusion model. On D3IL-aligning analytic
    average-velocity matching closes $84\,\%$ of the distance to the target against $14$, $10$ and
    $-2\,\%$ for the others, and under projection keeps nine contexts violation-free against the
    baseline's four at a third of its time. Instantaneous-velocity matching and the diffusion model
    are level on outcome there and on D3IL-avoiding, and differ in cost. Between the two
    average-velocity objectives the consistency-interpolated one is ahead by one episode in thirty on
    D3IL-avoiding and behind by most of the task on D3IL-aligning, where it does not improve with
    the budget; the analytic objective leads or holds in every environment. Where the demonstrations
    are simple, on UAV-corridor and UAV-s-curve, the objectives are equivalent or
    instantaneous-velocity matching leads, and the budget orders the configurations instead. One
    evaluation saturates D3IL-avoiding, D3IL-aligning needs twenty, and on the s-curve more
    evaluations lose flights.
  \item[RQ2 --- Constraints.] Endpoint projection guides the sample only where the budget and the
    threshold leave a sampling step before the terminal one, $n\sidx{guide} = \lceil \eta\nfe \rceil
    - 1$: at one evaluation never, at two at the default threshold not, and there it computes the same
    plan as per-step projection at $2.45$ times the cost. Where it has steps to guide it is the better
    projector on the constraint at a comparable price: on D3IL-aligning at twenty evaluations it keeps
    all ten contexts violation-free where per-step projection keeps nine, in seven of nine matched
    comparisons, and at ten evaluations at $37\,\%$ of per-step projection's time; on D3IL-avoiding at
    three evaluations it halves the time per control step of the two average-velocity models at
    $1.000$. It does not reach the goal in fewer control steps, and on UAV-corridor it trades: fewer
    and shallower violations at twenty evaluations, cheaper over the hump with one candidate and
    dearer under the tilt, on longer flights. It helps on the hard task that needs the budget,
    D3IL-aligning, and not on D3IL-avoiding, which is solved at a budget where it has no step to act.
  \item[RQ3 --- Transfer.] The answers hold where the demonstrations are of the same kind, and the
    plant and the demonstrations change the rest. Analytic average-velocity matching leads again on
    camera observations and contact, and the declared constraints transfer exactly from the
    manipulator to the quadrotor on UAV-pillars, where the two average-velocity models keep their
    order. The quadrotor adds what the manipulator could not show: an undeclared obstacle ends a
    flight, the controller decides whether a plan on a route of sharp turns is flown at all, and on
    demonstrations that carry one path forward from every position the objectives cannot be told
    apart. The projection scheme itself carries over to three dimensions and to constraints that leave
    the horizontal plane.
\end{description}
